\documentclass[11pt,a4paper]{article}

\usepackage{amsmath,amssymb,bm}
\usepackage[margin=25mm]{geometry}

\newcommand{\chip}{\chi^{+}}
\newcommand{\chim}{\chi^{-}}
\newcommand{\eps}{\epsilon}
\newcommand{\bbeta}{\beta}

\title{Hamiltonian Density in the $\chi^\pm$ Basis}
\author{}
\date{}

\begin{document}
\maketitle

\section{Definition of the $\chi^\pm$ modes}

We use the two local Majorana-like modes
\begin{align}
  \chi_j^\pm
  =
  \frac{1}{\sqrt{2}}
  \left(
    e^{\pm i\pi/4}c_j
    +
    e^{\mp i\pi/4}c_j^\dagger
  \right).
\end{align}
The inverse transformation is
\begin{align}
  c_j
  =
  \frac{1}{\sqrt{2}}
  \left(
    e^{-i\pi/4}\chi_j^+
    +
    e^{i\pi/4}\chi_j^-
  \right).
\end{align}
The anticommutation relations are
\begin{align}
  \{\chi_i^+,\chi_j^+\}=\delta_{ij},
  \qquad
  \{\chi_i^-,\chi_j^-\}=\delta_{ij},
  \qquad
  \{\chi_i^+,\chi_j^-\}=0.
\end{align}

\section{Useful bilinears}

The onsite number operator becomes
\begin{align}
  c_j^\dagger c_j
  =
  \frac{1}{2}
  +
  i\chi_j^+\chi_j^- .
\end{align}
For a nearest-neighbor bond $(j,j+1)$,
\begin{align}
  c_{j+1}c_j
  =
  \frac{1}{2}
  \left(
    i\chi_j^+\chi_{j+1}^+
    -
    i\chi_j^-\chi_{j+1}^-
    -
    \chi_j^-\chi_{j+1}^+
    -
    \chi_j^+\chi_{j+1}^-
  \right),
\end{align}
and
\begin{align}
  c_{j+1}^\dagger c_j
  =
  \frac{1}{2}
  \left(
    -\chi_j^+\chi_{j+1}^+
    -
    \chi_j^-\chi_{j+1}^-
    -
    i\chi_j^-\chi_{j+1}^+
    +
    i\chi_j^+\chi_{j+1}^-
  \right).
\end{align}
The corresponding expressions for the left bond $(j-1,j)$ are obtained by
replacing $(j,j+1)$ with $(j-1,j)$.

\section{Site-centered density used in the code}

The code defines a site density by averaging the left and right bonds around
site $j$.  In this convention, the same local factor
$\bbeta_{j+1/2}\equiv \bbeta(j+1/2)$ multiplies the averaged
nearest-neighbor operator at site $j$.
Therefore the site-centered densities are
\begin{align}
  H_j^+
  &=
  -\frac{i}{4\eps}
  \left(1+\bbeta_{j+1/2}\right)
  \left(
    \chi_{j-1}^+\chi_j^+
    +
    \chi_j^+\chi_{j+1}^+
  \right),
  \\
  H_j^-
  &=
  \frac{i}{4\eps}
  \left(-1+\bbeta_{j+1/2}\right)
  \left(
    \chi_{j-1}^-\chi_j^-
    +
    \chi_j^-\chi_{j+1}^-
  \right),
  \\
  H_j^{+-}
  &=
  -\frac{i p}{4\eps}
  \left(
    \chi_{j-1}^+\chi_j^-
    +
    \chi_{j-1}^-\chi_j^+
    +
    \chi_j^+\chi_{j+1}^-
    +
    \chi_j^-\chi_{j+1}^+
  \right)
  -
  \frac{p-\eps m}{\eps}\chi_j^+\chi_j^- .
\end{align}
The local Hamiltonian density is
\begin{align}
  H_j
  =
  H_j^+ + H_j^- + H_j^{+-}.
\end{align}
For open boundary conditions, the implementation sets the endpoint densities
to zero.

\section{Total Hamiltonian in the code convention}

The total Hamiltonian is the sum over site-centered densities,
\begin{align}
  H = \sum_j H_j .
\end{align}
Because each nearest-neighbor bond appears once from the density at its left
site and once from the density at its right site, the coefficient on the bond
$(a,a+1)$ contains the averaged beta factor
\begin{align}
  \bar{\bbeta}_{a+1/2}
  =
  \frac{\bbeta_{a+1/2}+\bbeta_{a+3/2}}{2}.
\end{align}
This is the same link average used in the BdG matrix construction.  Therefore
the summed Hamiltonian can be written as
\begin{align}
  H
  &=
  \sum_a
  \Bigg[
    -\frac{i}{2\eps}
    \left(1+\bar{\bbeta}_{a+1/2}\right)
    \chi_a^+\chi_{a+1}^+
    +
    \frac{i}{2\eps}
    \left(-1+\bar{\bbeta}_{a+1/2}\right)
    \chi_a^-\chi_{a+1}^-
    \nonumber\\
  &\hspace{35mm}
    -
    \frac{i p}{2\eps}
    \left(
      \chi_a^+\chi_{a+1}^-
      +
      \chi_a^-\chi_{a+1}^+
    \right)
  \Bigg]
  -
  \frac{p-\eps m}{\eps}
  \sum_j
  \chi_j^+\chi_j^- .
\end{align}

\section{Heisenberg equation of motion}

For either species $\sigma=\pm$, the Heisenberg equation is
\begin{align}
  i\frac{\partial}{\partial t}\chi_j^\sigma
  =
  [\chi_j^\sigma,H].
\end{align}
Using the anticommutation relations, a useful identity for any two
$\chi$ operators is
\begin{align}
  [\chi_a,\chi_b\chi_c]
  =
  \{\chi_a,\chi_b\}\chi_c
  -
  \chi_b\{\chi_a,\chi_c\}.
\end{align}
Equivalently, if the collective index $\alpha=(j,\sigma)$ is used and
$\{\chi_\alpha,\chi_\beta\}=\delta_{\alpha\beta}$, then
\begin{align}
  [\chi_\alpha,\chi_\beta\chi_\gamma]
  =
  \delta_{\alpha\beta}\chi_\gamma
  -
  \delta_{\alpha\gamma}\chi_\beta .
\end{align}

\subsection{Equations from the code Hamiltonian}

Using the summed Hamiltonian in the code convention, the equations of motion
are
\begin{align}
  i\partial_t \chi_j^+
  &=
  -\frac{i}{2\eps}
  \left[
    \left(1+\bar{\bbeta}_{j+1/2}\right)\chi_{j+1}^+
    -
    \left(1+\bar{\bbeta}_{j-1/2}\right)\chi_{j-1}^+
  \right]
  \nonumber\\
  &\quad
  -\frac{i p}{2\eps}
  \left(
    \chi_{j+1}^-
    -
    \chi_{j-1}^-
  \right)
  -
  \frac{p-\eps m}{\eps}\chi_j^- ,
  \\
  i\partial_t \chi_j^-
  &=
  \frac{i}{2\eps}
  \left[
    \left(-1+\bar{\bbeta}_{j+1/2}\right)\chi_{j+1}^-
    -
    \left(-1+\bar{\bbeta}_{j-1/2}\right)\chi_{j-1}^-
  \right]
  \nonumber\\
  &\quad
  -\frac{i p}{2\eps}
  \left(
    \chi_{j+1}^+
    -
    \chi_{j-1}^+
  \right)
  +
  \frac{p-\eps m}{\eps}\chi_j^+ .
\end{align}
Here
\begin{align}
  \bar{\bbeta}_{j+1/2}
  =
  \frac{\bbeta_{j+1/2}+\bbeta_{j+3/2}}{2},
  \qquad
  \bar{\bbeta}_{j-1/2}
  =
  \frac{\bbeta_{j-1/2}+\bbeta_{j+1/2}}{2}.
\end{align}
These follow directly from
$i\partial_t\chi_j^\pm=[\chi_j^\pm,H]$ and the bilinear commutator identity
above.

\end{document}
